% Volume III: Attention and Publicity
% Part VI. Privacy as Boundary and Capacity
% Chapter 31: Contextual Integrity
% Spine: proof-spines/ch031-spine.md

\chapter{Contextual Integrity}
\label{ch:ch031}

\textbf{Contextual integrity} represents an information flow as
\[
\tau=(s,r,\alpha,c,p),
\]
\Definition\ where \(s\) is sender, \(r\) recipient, \(\alpha\) information type, \(c\) context, and \(p\) transmission principle. On this account, many privacy violations are not disclosures of secrets at all, but unauthorized changes in \(\tau\): the same content crosses into a different role-structure, purpose, or principle of transmission.

This framework replaces the crude public/private binary with a more exact question: who may know what, in which role, in which context, and under which transmission constraints? Information appropriate in one setting may become harmful when moved into another, even when nothing about the underlying fact changes. Privacy protects intelligibility by preserving these context-sensitive norms of flow. The same norm-preserving restrictions become concealment when institutions invoke context to block adversarial scrutiny of actors exercising coercive or allocative power.
